% Part: incompleteness
% Chapter: theories-computability
% Section: oconsis-ext-of-q-undec

\documentclass[../../../include/open-logic-section]{subfiles}

\begin{document}

\olfileid{inc}{tcp}{oqn}

\olsection{$\Th{Q}$-এর $\omega$-সঙ্গতিপূর্ণ প্রসারণগুলি অনির্ণেয়}

\begin{explain}
$\Th{Q}$ যে c.e.-সম্পূর্ণ, তার প্রমাণে এই তথ্যটি ব্যবহৃত হয়েছিল যে
$\Th{Q}$-তে প্রমাণযোগ্য প্রতিটি বাক্য স্বাভাবিক সংখ্যাগুলির ক্ষেত্রে
``সত্য''। পরের সংজ্ঞা ও উপপাদ্যটি উপরের প্রমাণে ``সত্য''-এর ঠিক কোন
দিকগুলি দরকার হয়েছিল তা চিহ্নিত করে ফলটিকে শক্তিশালী করে। তবে এই উপপাদ্য
নিয়ে খুব বেশি সময় দিও না, কারণ শিগগিরই আমরা আরও শক্তিশালী ফল প্রমাণ করব।
মূলত ঐতিহাসিক কারণে এটি অন্তর্ভুক্ত করা হয়েছে: গ্যোডেলের মূল প্রবন্ধে
$\omega$-সঙ্গতির ধারণা ব্যবহৃত হয়েছিল, কিন্তু অল্পকাল পরেই
$\omega$-সঙ্গতির জায়গায় সাধারণ সঙ্গতি বসিয়ে তাঁর ফলটিকে শক্তিশালী করা হয়।
\end{explain}

\begin{defn}
\ollabel{thm:oconsis-q}
কোনো তত্ত্ব~$\Th{T}$-কে $\omega$-সঙ্গতিপূর্ণ বলা হয় যদি নিচের শর্তটি
পূরণ হয়: $\lexists[x][!A(x)]$ কোনো বাক্য এবং $\Th{T}$ যদি
$\lnot !A(\num 0)$, $\lnot !A(\num 1)$, $\lnot !A(\num 2)$,~\dots
প্রমাণ করে, তবে $\Th{T}$ $\lexists[x][!A(x)]$ প্রমাণ করে না।
\end{defn}

\begin{thm}
  ধরি $\Th{T}$ এমন যেকোনো $\omega$-সঙ্গতিপূর্ণ তত্ত্ব যা $\Th{Q}$-কে
  অন্তর্ভুক্ত করে। তাহলে $\Th{T}$ অনির্ণেয়।
\end{thm}

\begin{proof}
$\Th{T}$ যদি $\Th{Q}$-কে অন্তর্ভুক্ত করে, তবে $\Th{T}$ গণনসাধ্য অপেক্ষক
ও সম্বন্ধগুলিকে উপস্থাপন করে। আগের প্রমাণটিকে শুধু সামান্য বদলাতে হবে। আগের
মতোই, $x \in K$ হলে $\Th{T}$ সূত্র
$\lexists[s][!A_T(\num x,\num x,s)]$ প্রমাণ করে। বিপরীত দিকে, ধরি
$\Th{T}$ সূত্র $\lexists[s][!A_T(\num x, \num x, s)]$ প্রমাণ করে। তখন
$x$ অবশ্যই $K$-তে থাকবে: তা না হলে নিবেশ~$x$-এ যন্ত্র~$x$-এর কোনো থামা
গণনা নেই; আর $!A_T$ ক্লিনির $T$ সম্বন্ধকে উপস্থাপন করে বলে $\Th{T}$
$\lnot !A_T(\num x, \num x, \num 0)$,
$\lnot !A_T(\num x, \num x, \num 1)$,~\dots প্রমাণ করবে। এতে
$\Th{T}$ $\omega$-অসঙ্গতিপূর্ণ হয়ে যাবে।
\end{proof}

\end{document}
